\section*{Chapter \ref{ch:b3:homogeneous-catalysis} --- Homogeneous Catalysis in Industry}

\begin{solution}{exo:b3:homogeneous-catalysis:1}
$\ce{RhCl(PPh3)3}$: 16, Rh(I) $d^8$. $\ce{RhCl(PPh3)2}$: 14, Rh(I). $\ce{RhClH2(PPh3)2}$: 16, Rh(III) $d^6$.
Alkene complex: 18, Rh(III). Alkyl hydride: 16, Rh(III). Reductive elimination
returns to 14, Rh(I).
\end{solution}

\begin{solution}{exo:b3:homogeneous-catalysis:2}
Butanal, $\ce{CH3CH2CH2CHO}$ (linear), and 2-methylpropanal, $\ce{(CH3)2CHCHO}$ (branched).
\end{solution}

\begin{solution}{exo:b3:homogeneous-catalysis:3}
Diethyl cyclopent-3-ene-1,1-dicarboxylate, a five-membered ring, with loss of
ethene: the two terminal $\ce{CH2}$ groups leave as $\ce{C2H4}$ and the two internal CH carbons
join.
\end{solution}

\begin{solution}{exo:b3:homogeneous-catalysis:4}
Methyl (\textit{E})-3-phenylprop-2-enoate (methyl cinnamate): the aryl adds to the
terminal carbon, and $\beta$-hydride elimination, syn, from the more stable conformer
gives the \textit{E} alkene.
\end{solution}

\begin{solution}{exo:b3:homogeneous-catalysis:5}
The oxidative addition is slow and $\ce{[Rh(CO)2I2]-}$ is the resting state, so
$v = k[\mathrm{Rh}]_{\mathrm{tot}}[\ce{CH3I}]$. Methanol only regenerates methyl iodide by a fast reaction with HI;
the total iodide charged sets $[\ce{CH3I}]$, so the rate does not depend on the methanol
concentration until there is too little methanol left to convert HI back into
methyl iodide.
\end{solution}

\begin{solution}{exo:b3:homogeneous-catalysis:6}
Phenylboronic acid with 4-bromotoluene, or 4-methylphenylboronic acid with
bromobenzene. Both work; the first uses the cheapest boronic acid and a stable,
common aryl bromide.
\end{solution}

\begin{solution}{exo:b3:homogeneous-catalysis:7}
$C_2$-symmetric: isotactic (both sites select the same face). $C_s$-symmetric:
syndiotactic (mirror-image sites alternate). Unbridged $\ce{Cp2ZrCl2}$: atactic (achiral
sites).
\end{solution}

\begin{solution}{exo:b3:homogeneous-catalysis:8}
$\mathrm{TON} = 0.98/0.0010 = 980$; mean $\mathrm{TOF} = 980/\qty{2.0}{h} = \qty{490}{h^{-1}}$.
\end{solution}

\begin{solution}{exo:b3:homogeneous-catalysis:9}
The repeating unit is $\ce{-[CH=CH-C5H8]-}$ with the two CH groups on carbons 1 and 3 of a
cyclopentane ring (poly(1,3-cyclopentylenevinylene)). The release of the
strain of the bicyclic ring makes $\Delta_rH^\circ$ strongly negative, which outweighs the
loss of translational entropy; no gas is formed, so the driving force is
enthalpic.
\end{solution}

\begin{solution}{exo:b3:homogeneous-catalysis:10}
$\bar n_n = 1/(1 - p) = 2000$; $M_w/M_n = 1 + p = 1.9995 \approx 2.0$. A \omterm{def:b3:homogeneous-catalysis:ziegler-natta}{Ziegler--Natta catalyst} has
several kinds of sites, each with its own $p$; the sum of several Schulz--Flory
distributions with different means is broader (dispersity often 4 to 8).
\end{solution}

\begin{solution}{exo:b3:homogeneous-catalysis:11}
At low pressure $K_1K_2[\ce{H2}] \ll [\mathrm L] + K_1$: $v \approx kK_1K_2[\mathrm{Rh}]_{\mathrm{tot}}[\text{alkene}][\ce{H2}]/([\mathrm L] + K_1)$,
first order in $\ce{H2}$. At high pressure the term $K_1K_2[\ce{H2}]$ dominates: $v \to k[\mathrm{Rh}]_{\mathrm{tot}}[\text{alkene}]$,
independent of $\ce{H2}$. $[\mathrm L]$ appears only in the denominator: added phosphine lowers $v$
by pushing rhodium back into the saturated precatalyst.
\end{solution}

\begin{solution}{exo:b3:homogeneous-catalysis:12}
The insertion that puts rhodium on the internal carbon builds a secondary alkyl
next to a crowded metal; bulky phosphines (and an excess of them, which keeps two
on the metal) make that transition state worse than the one giving the primary
alkyl, which leads to the linear aldehyde. Linear butanal is the one converted by
aldol condensation and hydrogenation into the branched $C_8$ alcohol used in
plasticisers; 2-methylpropanal is a lower-value by-product.
\end{solution}

\begin{solution}{pb:b3:homogeneous-catalysis:1}
\textbf{1.} Ir(I), $d^8$: $9 + 4 + 2 + 1 = 16$ electrons (neutral counting, with the charge).

\textbf{2.} $\ce{[CH3Ir(CO)2I3]-}$, 18 electrons, Ir(III).

\textbf{3.} $\ce{[CH3Ir(CO)3I2]}$, 18 electrons, Ir(III).

\textbf{4.} $\ce{[CH3COIr(CO)2I2]}$, 16 electrons, Ir(III).

\textbf{5.} $\ce{[CH3COIr(CO)2I3]-}$ (18) eliminates $\ce{CH3COI}$, regenerating $\ce{[Ir(CO)2I2]-}$.

\textbf{6.} $\ce{CH3COI + H2O -> CH3COOH + HI}$; $\ce{CH3OH + HI -> CH3I + H2O}$.

\textbf{7.} $\ce{CH3OH + CO -> CH3COOH}$.

\textbf{8.} With iridium the migratory insertion is slow; the oxidative addition of
methyl iodide, rate-determining with rhodium, is much faster on the more
electron-rich iridium.

\textbf{9.} They remove one iodide from $\ce{[CH3Ir(CO)2I3]-}$, giving the neutral tricarbonyl, in which
the metal back-donates less and the methyl migrates faster onto a more
electrophilic CO.

\textbf{10.} An inverse dependence on the free iodide concentration: iodide pushes
the pre-equilibrium back to the slow anionic complex.

\textbf{11.} Methanol enters only after the rate-determining step, converted into
methyl iodide by a fast reaction with HI.

\textbf{12.} $M(\ce{CH3COOH}) = \qty{60.1}{g/mol}$: $n = \qty{5.0e11}{g}/\qty{60.1}{g/mol} = \qty{8.3e9}{mol}$ per year.

\textbf{13.} $\qty{8.32e9}{mol}/\qty{8000}{h} = \qty{1.04e6}{mol/h}$.

\textbf{14.} $\qty{2.0e5}{L} \times \qty{5.0e-3}{mol/L} = \qty{1.0e3}{mol}$ of iridium, $\qty{1.0e3}{mol} \times \qty{192.2}{g/mol} \approx \qty{190}{kg}$.

\textbf{15.} $\mathrm{TOF} = \qty{1.04e6}{mol/h}/\qty{1.0e3}{mol} = \qty{1.0e3}{h^{-1}}$, about \qty{0.29}{s^{-1}}.

\textbf{16.} $\mathrm{TON} = \qty{8.3e9}{mol}/\qty{1.0e3}{mol} = \num{8.3e6}$ per year.

\textbf{17.} $\qty{5.0e8}{kg}/(\qty{8000}{h} \times \qty{200}{m^3}) \approx \qty{310}{kg.m^{-3}.h^{-1}}$.

\textbf{18.} $\ce{CO + H2O -> CO2 + H2}$.

\textbf{19.} $1/0.94 = \qty{1.064}{mol}$ of CO per mole of acetic acid.

\textbf{20.} $1.064 - 1 = \qty{0.064}{mol}$ of \ce{CO2} per mole of acetic acid.

\textbf{21.} One tonne is $\qty{1.0e6}{g}/\qty{60.1}{g/mol} = \qty{1.66e4}{mol}$; $\qty{1.66e4}{mol} \times 0.0638 \times \qty{44.0}{g/mol} \approx \qty{47}{kg}$ of \ce{CO2}.

\textbf{22.} Water is a reactant of the shift: less water slows it and saves the energy
of drying the acid. Water is still needed to hydrolyse acetyl iodide and to keep the
catalyst active and dissolved; with rhodium, low water precipitates rhodium
iodide.

\textbf{23.} Hydrogen: it builds up in the gas, lowers the partial pressure of CO
(so gas must be vented, losing CO), and can hydrogenate intermediates into
by-products such as propionic acid.

\textbf{24.} The iridium catalyst turns over about $\qty{1.0e3}{h^{-1}}$: each iridium atom makes a
thousand molecules of acetic acid an hour, some eight million a year.
\end{solution}
